% ch-ref-boundary-output.tex -- reference: the Boundary and Output blocks
\chapter{Boundary conditions and output}
\label{ch:ref-bcout}

Two blocks are documented here. \key{Boundary} decides what happens at the
edges of the domain, which is one of the choices that most strongly changes the
answer; \key{Output} decides what is written to disk and how often, which
decides whether you can tell what happened. Both are required and both appear
exactly once.

\section{\key{Boundary}}
\label{sec:ref-bcout:boundary}

\key{Boundary} carries no keys of its own. It contains one sub-block per
domain edge, each with exactly one key, \key{Type}. \key{XminBC} and
\key{XmaxBC} are required; \key{YminBC} and \key{YmaxBC} are optional and are
only meaningful when \key{Grid.Dimension} is \val{"xyh"}.

\begin{keytable}{The four sub-blocks of \key{Boundary}, each with the single key \key{Type}\label{tab:boundary-keys}}\key{XminBC.Type} & string & \val{"Symmetry"} & Condition on the $x=\mathrm{Xmin}$ edge: \val{"Symmetry"}, \val{"Periodic"}, \val{"Pinning"}, \val{"FreshInlet"} or \val{"Outlet"}. The block is required. \\
\key{XmaxBC.Type} & string & \val{"Symmetry"} & Condition on the $x=\mathrm{Xmax}$ edge, same set of values. The block is required. \\
\key{YminBC.Type} & string & \val{"Symmetry"} & Condition on the $y=\mathrm{Ymin}$ edge, same set of values. The block is optional and is ignored in 1-D. \\
\key{YmaxBC.Type} & string & \val{"Symmetry"} & Condition on the $y=\mathrm{Ymax}$ edge, same set of values. The block is optional and is ignored in 1-D. \\
\end{keytable}

\begin{note}[Two conditions per edge, because the equation is fourth order]
  The height equation contains a fourth spatial derivative, so each boundary
  needs \emph{two} conditions, not one. That is why you choose a named type
  rather than writing a value: each name stands for a pair of conditions, and
  the ghost layer is two cells deep so that the pair can be imposed. The
  concentration equation is second order and takes one condition from the same
  name.
\end{note}

\subsection{What each type means}

\paragraph{\val{"Symmetry"}} A plane of symmetry: no flux crosses the
boundary, and both the shape and the concentration are mirrored through it,
$\partial h/\partial n = 0$ and $\partial^3 h/\partial n^3 = 0$: the two ghost
layers are the exact mirror image of the two layers inside, in 1-D as in 2-D
(\cref{sec:time:apply}). Use it to put
the center of a drop or a film on the boundary and halve the cost, or as a
closed wall. It is the default. It is the wrong choice whenever liquid should
be able to leave: a drop spreading toward a \val{"Symmetry"} edge behaves as
if a wall stood there, and gives a different answer from \val{"Outlet"}.

\paragraph{\val{"Periodic"}} The edge is joined to the opposite one; the ghost
layers are copied across. Use it for a film on a periodically structured
substrate, or to compute one wavelength of a pattern that continues for ever.
It is the wrong choice when the physics has a net drift across the domain ---
with in-plane gravity or a source, periodicity recycles what flows out of one
side back into the other. With the centrifugal term it is refused: the arm of
the force, the distance from the rotation axis, jumps by the length of the
domain across the seam, so a rotating film is not periodic
(\cref{sec:terms:centrifugal}; a run with a manufactured solution is
exempt).

\paragraph{\val{"Pinning"}} The contact line is held at the edge: the film
thickness at the boundary is pinned to zero and the in-plane velocity there is
zero, while the concentration takes a zero-gradient condition. Because the thickness is held at zero, \textbf{the initial film must reach
zero at a pinned edge too}, and the solver checks it before the first step: it
extrapolates the film from the first two cells to the wall and reports the
value in the run log,
\begin{plain}
Pinning Xmax: the initial film reaches 0.204 um at the wall (2.0 % of its 9.995 um maximum) -- ok.
\end{plain}
and stops the run, with that line and what to do, if more than a tenth of the
thickest film is left at the wall. A film of full thickness against a pinned
edge is a step inside one cell, which the first steps cannot get across: the
run would otherwise end in stagnating Newton iterations and a time step cut
below \key{MinDt}, with nothing in the messages pointing at the wall.
\key{Initial.\bk SurfaceShape.\bk EdgeBump} at that \key{Location} is the shape
made for this, and it needs a small \key{PrecursorRatio}, because the film is
never thinner than the precursor anywhere else. In 2-D every row along the edge
is checked and the worst one is reported. The check is of the initial
condition: a restart (\key{Restart.Enable}) continues from the film of its
field file, whose edge may have evolved --- pushed against the wall and dried,
the film there can pass a tenth of the thickest --- and is taken as it is, the
line written as a warning when it is past that limit
(\cref{sec:ref-numerics:restart}). It represents a
liquid edge arrested by a step, an edge bead or a contamination line, and it
lets you model a film whose edge does not move \emph{without} paying for a
resolved contact line. It is the wrong choice if the question is how fast the
edge moves --- for that you need the disjoining pressure and a mesh that
resolves the transition region (\cref{sec:terms:disjoining}). It is also the
type that \key{Initial.SurfaceShape.EdgeBump} requires at its own
\key{Location}.

\paragraph{\val{"FreshInlet"}} An inflow boundary for liquid of constant
composition: the concentration at the edge is held at its initial value, the
thickness at the edge face is taken from the current film, and
$\partial^3 h/\partial n^3 = 0$. Use it at the upstream end of a coating. The
composition of what flows in is prescribed; how much flows in is decided by
the film.

\begin{note}[\val{"FreshInlet"} has no keys, so the inlet state comes from \key{Initial}]
  There is no key anywhere in the file for the inlet thickness or the inlet
  concentration. The prescribed values are the average of the two innermost
  points at that edge --- taken from the \emph{current} field for the height
  equation, and from the field snapshot saved when the boundary condition was
  built (\code{\_field\_ini}) for the concentration equation. The inlet
  therefore follows from your \key{Initial} block, and the concentration side
  in particular is frozen at its initial value for the whole run. To change
  what flows in, change the initial profile near that edge.
\end{note}

\paragraph{\val{"Outlet"}} Free outflow: $\partial^2 h/\partial n^2 = 0$ and
$\partial^3 h/\partial n^3 = 0$, so liquid leaves without the boundary imposing
a shape on it. Use it where the domain was simply cut for cost and you have no
information about what lies beyond. The volume that leaves, and the solute
with it, are summed from the face fluxes of each step and reported in
\file{History.dat} (\code{OutflowedVol}, \code{OutflowedResin}; what comes
in through a \val{"FreshInlet"} counts there negatively), so that
\code{VolErr} and \code{ResinVolErr} stay at rounding
(\cref{sec:output:history}); that is how you check that an outlet is doing
what you think. It is the wrong choice where a real wall or a real symmetry
plane exists, because it will let liquid drain away that ought to have piled
up.

\subsection{Combinations and edge cases}

Set \val{"Periodic"} on \emph{both} edges of a pair or on neither: a single
periodic edge is not a periodic domain. The input consistency check
(\code{CheckParameters::Check}) refuses a one-sided \val{"Periodic"}, in either
direction and in 1-D as well as 2-D, before anything is built, with the message
\code{Periodic must be set on both edges of a pair or on neither}. It refuses
a \val{"Periodic"} edge with the centrifugal term as well (a rotating film is
not periodic; \code{a rotating film is not periodic} in the message), unless
\key{Numerics.ManufacturedSolution} is on.

The \val{"CZ"} model implements two kinds of edge only, \val{"Symmetry"} and
\val{"Periodic"}. The consistency check refuses \val{"Pinning"},
\val{"FreshInlet"} and \val{"Outlet"} for \val{"CZ"}, naming the edge.

In two dimensions the corners are handled for every pair of adjacent types
except \val{"Pinning"} combined with \val{"FreshInlet"} or with \val{"Outlet"}.
Those four corner cases --- both types in either order --- are \emph{not
implemented, and they abort the run}: the message \code{not yet implemented}
goes to \file{FatalErrors.txt} and to standard error, and nothing is written to
\file{Warnings.txt} (\file{model/\bk BoundaryCondition.cpp}). If you need a
pinned edge next to an inlet or an outlet, you will have to implement that
corner (\cref{sec:extending:other}).

\begin{warn}[The substrate is mirrored at a wall]
  Whatever you specify for the film and the concentration, the
  \key{Substrate} shape is continued into the ghost layers as its mirror image
  at every edge except a \val{"Periodic"} one, where it is continued from the
  other end of the domain. A substrate topography that is not symmetric about
  a wall --- a plane tilted by \key{FlatPlane}, or a bump cut in half by the
  domain edge --- is therefore met in the ghost layer by its mirror image, not
  by its own continuation. Place a wall where the substrate is flat, or where
  its slope is genuinely zero, whenever you can.
\end{warn}

\section{\key{Output}}
\label{sec:ref-bcout:output}

\key{Output} decides what is written into the result directory named by
\key{JobName}, and how often. It has two keys of its own; \key{FieldData},
\key{GeneralHistory} and \key{CalcInfo} are required sub-blocks, and
\key{ProbeList} is optional.

\begin{keytable}{Keys of \key{Output}\label{tab:output-keys}}\key{Overwrite} & bool & \val{false} & Allow the run to proceed when the result directory already exists, replacing the files it writes again (\cref{sec:output:dir}). When false, an existing directory aborts the run before anything is written. \\
\key{Dimensioned} & bool & \val{true} & Write results in dimensional units instead of the solver's non-dimensional ones. Time is in \si{\second} and velocities in \si{\metre\per\second} everywhere; the length units differ between writers, see below. It also decides how a restart file is interpreted; see \cref{sec:ref-numerics:restart}. \\
\end{keytable}

\begin{warn}[\key{Dimensioned} does not put every file in the same length unit]
  Heights are in \si{\micro\metre} and in-plane coordinates in
  \si{\milli\metre} in everything that reports a position: \file{field*.dat},
  \file{field*.vts} and \file{ProbeNN.dat}. \file{History.dat} is the
  exception and takes the raw scales, so its lengths are in \si{\metre} and
  its volumes in \si{\metre\cubed} (in 1-D, per meter of width,
  \si{\metre\cubed\per\metre}); its volumes would otherwise have to be
  in a mixed \si{\milli\metre\squared\micro\metre}, which is worse than the
  difference between two files.

  Every one of these files states its units in its header, so there is no need
  to remember the table: read the first line. Do check it before comparing a
  number taken from one file with a number taken from another.
\end{warn}

\subsection{What gets written}

\begin{description}
  \item[\file{field00000.vts}, \dots] The full field as a binary VTK
    structured grid. Beside it sits \file{<JobName>.pvd}, the time-series
    index that makes ParaView treat the sequence as one animation. Controlled
    by \key{FieldData.VTKFile}.
  \item[\file{field00000.dat}, \dots] The same field as tab-separated ASCII.
    This is also the ASCII restart format, so any field dump can be used to
    continue a run, and it is hand-editable, which is how an analytically
    constructed profile is injected. Controlled by
    \key{FieldData.NeutralFile}.
  \item[\file{History.dat}] Integral quantities and extrema over time:
    the columns after the time are listed in full in
    \cref{tab:history-columns}. This is the first file to look at after a run.
  \item[\file{CalcInfo.dat}] Solver diagnostics, one row per output interval: Newton iterations and
    residual, linear-solver iterations and residual, the time step, and the
    relative change rate that the steady-state test uses. Look here when a run
    is slow or stops early.
  \item[\file{Energy.dat}] With \key{Output.Energy}, the free energy of
    the film, its dissipation and their balance, with every row of
    \file{History.dat} (\cref{sec:output:energy}).
  \item[\file{ProbeNN.dat}] One time series per probe, if any probe was
    declared, with \code{NN} a zero-padded two-digit index counting from
    \code{00} in input-file order (\cref{sec:output:probes}).
  \item[\file{JobInfo.txt}, \file{Warnings.txt}, \file{FatalErrors.txt}] The
    text files. \file{JobInfo.txt} holds the program, machine and timing
    information and a few figures of the run (the largest Courant numbers,
    the solute moved past $c = 1$ and into cells below $c = 0$, \cref{sec:output:jobinfo}); what the solver built from your input --- the scales and
    the dimensionless groups it derived --- and how the run progressed go to
    standard output and to no file at all (\cref{sec:output:running}).
    \file{Warnings.txt} holds every warning the run raised and
    \file{FatalErrors.txt} the reason it aborted, the latter duplicated to
    standard error. Neither carries a banner, so each is empty exactly when
    nothing was raised; see \cref{sec:output:warnings}.
\end{description}

\subsection{How the intervals interact}
\label{sec:output:intervals}

All four \key{TimeInterval} keys are in \textbf{seconds} and are divided by the
time scale $T_0$ internally. Each of them also acts as a cap on the time step:
no step is longer than the interval, so no output is skipped, and each output
is written at the step whose time is nearest the scheduled one. The smallest interval in the file is
therefore an upper bound on $\Delta t$ for the whole run, alongside the bound
from \key{Numerics.TimeStepControl.MaxDtFactor}
(\cref{sec:ref-numerics:tsc}). Setting a very short history interval on a case
with an adaptive step is thus not free: it can hold the step below what the
physics requires.

\key{Numerics.InitialDt} is clamped down to the smallest output interval if it
is larger, so a first step longer than the first output never happens.

\textbf{The final state is always written}, whatever the intervals say. Each
kind of output is written before the step that follows it, so the state the run
finishes on --- the one at \key{EndTime}, which is the answer the run was asked
for --- would otherwise survive only when an interval happened to divide the
end time exactly. It is written once more when the loop ends, and only if that
time is not already in the file, so a schedule that does land on the end does
not produce it twice. A run stopped by a divergence is the exception: what it
holds is not a solution, so nothing further is written. The time stamps in the
\file{.pvd} are the times actually reached, not the scheduled ones, which is
why they carry the last few digits of the step history rather than round
numbers.

\subsection{\key{Output.FieldData}}

\begin{keytable}{Keys of \key{Output.FieldData}\label{tab:output-fielddata-keys}}\key{TimeInterval} & double & \val{0.01} & Interval between full field dumps, in seconds, positive. Also an upper bound on $\Delta t$. Not to be given with \key{Times} or a \key{LogSchedule}. \\
\key{NeutralFile} & bool & \val{true} & Write the ASCII neutral \file{.dat} field files, which double as the ASCII restart format. \\
\key{VTKFile} & bool & \val{true} & Write the binary \file{.vts} files and maintain the \file{.pvd} index for ParaView. \\
\key{WritePotential} & bool & \val{false} & Write the driving potential $\varphi$ as one more column \code{PHI} of \file{field*.dat} and one more array \code{Potential} of the \file{.vts} (\cref{sec:output:potential}). Well-mixed model only. \\
\key{Times} & array & \val{[]} & The times to write the field at instead of every \key{TimeInterval}, in seconds: ascending, positive. Those after \key{Numerics.EndTime} are dropped; at least one must remain, and an empty list is refused. \\
\key{ExactTimes} & bool & \val{false} & With \key{Times}: land a step on each listed time (\cref{sec:ref-output:listed}). Refused without \key{Times}. \\
\end{keytable}

\subsubsection{Output at listed times}
\label{sec:ref-output:listed}

A field written every \key{TimeInterval} holds the time step to that interval
for the whole run, so an output schedule that must resolve both the first
microseconds and the last seconds of a run --- a film levelling over seven
decades of time, say --- would hold every step to the shortest interval.
\key{Times}, or the sub-block \key{LogSchedule} below, lists the times
instead, and \textbf{a listed time caps only the step that would pass it}:
\begin{description}
  \item[\key{ExactTimes} \val{false}] the field is written at the step closest
    to each listed time, as with an interval, and the step is held to the gap
    between the two listed times it falls between --- as an interval holds it
    to the interval, but a gap that widens lets it lengthen. A step that
    passes several listed times writes once.
  \item[\key{ExactTimes} \val{true}] a step that would pass the next listed
    time is shortened to land on it, to rounding, and the field is written
    there. The step after it is the one the run was on before the
    shortening: a listed time is not a reason for the step control to start
    again from a short step. The shortened step is the step taken all the
    same --- the error estimate of the next step differences over it --- and
    when the step control asks for less than it, or the step was retried
    shorter, the run goes on from what the step control asks for. A
    remainder below $10^{-9}$ of the time is rounding and is not stepped to.
\end{description}
The state at the start is written first and the final state last, as with
an interval. A run continued from a \file{field*.state} goes on with the
listed time it was waiting for; continued from a field file alone, with the
first listed time after the restart time. \textbf{The other outputs still cap
the step at their intervals} (\cref{sec:output:intervals}): with listed
times, give \key{Output.GeneralHistory}, \key{Output.CalcInfo} and
\key{Output.ProbeList} intervals as long as the run allows.
\file{leveling\_1d\_logout} writes 141 times over seven decades, from
$10^{-7}$ to \SI{1}{\second}, in 403 steps; held to its shortest gap, the
same run would take at least $10^7$.

\begin{keytable}{Keys of \key{Output.FieldData.LogSchedule} (optional; instead of \key{TimeInterval} and \key{Times})\label{tab:output-logschedule-keys}}\key{Start} & double & \reqd & The first listed time, in seconds, positive. \\
\key{End} & double & \reqd & The last, in seconds, after \key{Start}; the listed times after \key{Numerics.EndTime} are dropped. \\
\key{PerDecade} & int & \reqd & Listed times per decade, at least 1: $\key{Start}\cdot 10^{i/\key{PerDecade}}$, $i = 0, 1, \dots$, up to \key{End}. \\
\key{ExactTimes} & bool & \val{false} & As \key{Output.FieldData.ExactTimes}. \\
\end{keytable}

\subsection{\key{Output.GeneralHistory}}

\begin{keytable}{Keys of \key{Output.GeneralHistory}\label{tab:output-history-keys}}\key{TimeInterval} & double & \val{0.1} & Interval of \file{History.dat}, in seconds, positive. Also an upper bound on $\Delta t$. \\
\end{keytable}

\subsection{\key{Output.CalcInfo}}

\begin{keytable}{Keys of \key{Output.CalcInfo}\label{tab:output-calcinfo-keys}}\key{TimeInterval} & double & \val{0.1} & Interval of \file{CalcInfo.dat}, in seconds, positive. Also an upper bound on $\Delta t$. \\
\end{keytable}

\subsection{\key{Output.ProbeList}}

A probe is a fixed point at which the solution is written on its own schedule,
much more often than a full field dump can afford to be. The block is
optional; left out, it reads as an empty one, so there are no probes and
\key{TimeInterval} is still \val{0.1}\,s, which still caps $\Delta t$.

\begin{keytable}{Keys of \key{Output.ProbeList}\label{tab:output-probelist-keys}}\key{TimeInterval} & double & \val{0.1} & Interval of the per-probe time series \file{ProbeNN.dat}, in seconds, positive. Read, and used to cap $\Delta t$, even when the list contains no probe at all. \\
\end{keytable}

\subsection{\key{Output.ProbeList.Probe}}

This block is \textbf{repeatable}: write one per probe, in the order you want
the files numbered. Write each one as an array-of-tables header,
\code{[[Output.ProbeList.Probe]]} (\cref{sec:inputformat:repeat}).

\begin{keytable}{Keys of \key{Output.ProbeList.Probe}, repeatable\label{tab:output-probe-keys}}\key{X} & double & \reqd & $x$ coordinate of the probe, in meters, any value. The nearest grid point within one cell is used; if there is none, the probe is silently dropped. \\
\key{Y} & double & \val{0} & $y$ coordinate, in meters. Ignored by the 1-D well-mixed model. \\
\key{Z} & double & \val{0} & Height above the substrate, in meters. Used only by the \val{"CZ"} model, which turns it into the $\zeta$ of its grid with the film thickness the probe's column starts with; the probe stays at that $\zeta$ as the film changes. \\
\end{keytable}

\begin{note}[Where a probe lands]
  No case in \file{cases/} declares a \key{Probe}; \code{ctest} checks them
  (\code{features}, \cref{sec:verification:features}). The header of
  \file{ProbeNN.dat} gives the point actually used, in millimeters when the
  output is dimensioned (and, for \val{"CZ"}, its $\zeta$).
\end{note}

\subsection{\key{Output.Energy}}

\begin{keytable}{Keys of \key{Output.Energy}\label{tab:output-energy-keys}}\key{Enable} & bool & \val{false} & Write \file{Energy.dat}, the free energy of the film, its dissipation and their balance, with every row of \file{History.dat} (\cref{sec:output:energy}). Well-mixed model only; refused with \val{"CZ"}. \\
\end{keytable}
